\section{Describing sparse pairs: validation and geometry}
\label{geom:section}

We now describe one side of a pair having unusually few incidences. The
decoder knows domains containing the two supports, but not the supports
themselves. The description will give a moderately enlarged set capturing
a fixed fraction of the smaller support. A separate validation step turns
this into almost complete capture on both sides. Dimensions at least four
are reduced by projection; the remaining geometric input concerns lines
and planes in projective dimension three. Section~\ref{learn:section}
supplies the sampling argument that completes the low-dimensional case.

Throughout this section, the parameters are those in
\eqref{base:parameters}; in particular, $q$ is prime and $P=o(\sigma)$,
where $\sigma=\log q$. All constants are uniform in the allowed parameters
and may depend on the fixed dimension bound, on $\eta$, and on fixed
constants in the hypotheses. The estimates are for sufficiently large
$q$. The reductions below change such constants only a bounded number of
times. This is harmless: $\log\sigma=o(K_*)$, so adding
$O(\log\sigma)$ to a deficiency bounded by a fixed multiple of $K_*$
still gives a bound of the same form.

\subsection{The description and its validation}

We use independent public random tables, each addressed separately. Given
a decoder-known finite set, a table provides independent rows of uniform
samples from that set. The encoder may test a row against hidden supports
and transmit the index of its first accepted row, up to a deterministic
cutoff. The decoder reconstructs the row from its index, without knowing
the acceptance predicate. Integers $m\geq 0$ are encoded with
$O(1+\log(m+1))$ bits. Support sizes, row lengths, orientations, and mode
choices are sent explicitly when needed. Every invocation below uses
tables independent of all data making that invocation ready.

\begin{lemma}[Sparse-pair description]\label{geom:description}
Let $2\leq j\leq d$. Let $S,T$ be nonempty sets in opposite point
systems of $\PG(j,q)$, contained in decoder-known sets $U_S,U_T$.
Suppose, for a fixed constant $C_b$, that
\begin{equation}\label{geom:sparse-hypotheses}
 |S||T|\geq q^{j+1}e^{-b},\qquad
 \frac{e(S,T)}{|S||T|}\leq\frac{\tau}{q},\qquad
 0\leq b\leq C_bK_*,\quad
 0<\tau\leq C_b\sigma^{-100\beta}.
\end{equation}
Write
\[
 d_S=\log\frac{|U_S|}{|S|},\qquad
 d_T=\log\frac{|U_T|}{|T|},
\]
and orient the systems so that $n=|S|\leq |T|$. There is a public-table
description, of length at most
\[
 CqP(d_S+d_T+P),
\]
which, except with probability $Ce^{-cq}$, produces a decoder-known set
$W\subseteq U_S$ satisfying
\begin{equation}\label{geom:description-output}
 |W|\leq ne^{CP},\qquad |W\cap S|\geq cn.
\end{equation}
The probability bound holds conditional on any earlier data independent
of the tables used by this invocation.
\end{lemma}

The proof is organized as dimension induction. We first prove validation
without assuming this lemma, then establish the induction step and the
geometric preparation in dimensions two and three. The final sampling
step is Proposition~\ref{learn:short-row}; the induction is closed
explicitly at the end of Section~\ref{learn:section}.

The sparse-incidence consequence of Lemma~\ref{base:incidence} gives
\begin{equation}\label{geom:product-upper}
 |S||T|\leq Cq^{j+1},\qquad
 n\leq Cq^{(j+1)/2},\qquad n\geq cq e^{-b}=q^{1-o(1)}.
\end{equation}
The last inequality uses $|T|\leq Q_j\leq Cq^j$. These bounds will be
used even when the product lower bound is weak by a factor $e^b$.
Figure~\ref{geom:validation-figure} separates the initial description
from the two fresh validation steps.

\begin{lemma}[Fresh validation]\label{geom:validation}
Under \eqref{geom:sparse-hypotheses}, suppose a decoder-known
$W\subseteq U_S$ satisfies \eqref{geom:description-output}. Further
independent tables and a message of length
$CqP(d_S+d_T+1)$ produce caps $U_T'\subseteq U_T$ and
$U_S'\subseteq U_S$, with additional failure probability $Ce^{-cq}$,
such that
\begin{equation}\label{geom:validation-output}
 \begin{aligned}
 |U_T'|&\leq \frac{Cq^{j+1}}{|S|},&
 |U_S'|&\leq \frac{Cq^{j+1}}{|T|},\\
 |U_T'\cap T|&\geq .99|T|,&
 |U_S'\cap S|&\geq .99|S|.
 \end{aligned}
\end{equation}
Each cap is its input domain intersected with an ambient test: a candidate
passes if it is incident with fewer than $h/(5q)$ entries of the
transmitted row of length $h$ in the opposite system.
\end{lemma}

\begin{proof}
Put $X=W\cap S$, so $|X|\geq c|S|$, and first consider an ideal row
of $h=\lceil C_0q(d_T+1)\rceil$ independent uniform points of $X$.
For an opposite point $y$, write
$p_X(y)=|X\cap N(y)|/|X|$, where $N(y)$ is its incident neighborhood.
The variance identity in Lemma~\ref{base:incidence} yields
\[
 \sum_y\bigl(p_X(y)-p_j\bigr)^2
 \leq \frac{q^{j-1}}{|X|},
 \qquad p_j=\frac{Q_{j-1}}{Q_j}.
\]
Since $qp_j=1-Q_j^{-1}$, there are at most
$Cq^{j+1}/|X|\leq Cq^{j+1}/|S|$ ambient points with
$p_X(y)<1/(2q)$. Every other point passes the row test with probability
at most $e^{-c_0h/q}$, by a lower binomial tail. The expected number of
such passing points in $U_T$ is at most
\[
 |U_T|e^{-c_0h/q}
 =|T|e^{d_T-c_0h/q}.
\]
Choosing $C_0$ large and then the cap-size constant large makes the
probability of violating the first size bound less than $.04$, by
Markov's inequality and $|T|\leq Cq^{j+1}/|S|$.

For capture, if $H_y$ denotes the row's hit count at $y$, then
\begin{equation}\label{geom:pointwise-iid}
 \Prob\{H_y\geq h/(5q)\}
 \leq 5q p_X(y).
\end{equation}
Averaging over $y\in T$ gives expected excluded fraction at most
\[
 \frac{5q e(X,T)}{|X||T|}
 \leq C\tau.
\]
Thus the probability of losing more than $.01|T|$ is $O(\tau)=o(1)$.
For large $q$, both desired properties hold for an ideal row with
probability at least $.9$.

To implement the ideal experiment, propose independent rows from the
known set $W$. Accept only rows whose entries all lie in $X$ and whose
ambient test has the two properties just proved. Conditional on all
entries lying in $X$, a proposal has exactly the ideal law. Moreover,
\[
 \frac{|W|}{|X|}\leq e^{C_1P}.
\]
A proposal is therefore accepted with probability at least
$.9e^{-C_1hP}$. Search at most
$\lceil e^{C_2hP}\rceil$ rows, with $C_2>C_1$ sufficiently large.
The failure probability is at most $e^{-cq}$, and the index costs
$O(hP)=O(qP(d_T+1))$ bits. The decoder intersects the resulting ambient
test with $U_T$, obtaining $U_T'$.

For the second test, the hidden source is $X_T=U_T'\cap T$ and its
proposal domain is $U_T'$. At readiness, the first successful test gives
\[
 \frac{|U_T'|}{|X_T|}
 \leq \frac{Cq^{j+1}}{|S||T|}
 \leq Ce^b\leq e^{CP}.
\]
Repeat the argument with $h=\lceil C_0q(d_S+1)\rceil$ and a fresh
table. The tested domain on the other side is the original $U_S$, not
$W$. Conditional on the first cap, its source is fixed, has at least
$.99|T|$ points, and satisfies
$e(S,X_T)/(|S||X_T|)\leq\tau/(.99q)$. Thus the same ideal-row
acceptance bound holds uniformly over every successful first cap.
Consequently this second test captures $.99$ of the original $S$.
The two costs and failure bounds give the conclusion; all finite
diagnostics cost $O(\sigma)$ and fit within the stated budget.
\end{proof}

\begin{figure}[t]
\centering
\begin{tikzpicture}[
  validationstage/.style={draw=black!55,align=center,text width=4.05cm,
    minimum height=3cm,inner sep=6pt,font=\small},
  flow/.style={-{Stealth[length=2mm]},thick},
  node distance=6mm]
\node[validationstage] (predict) {
  \textbf{Description}\\[5pt]
  Input domains $U_S,U_T$\\[3pt]
  Output $W\subseteq U_S$\\
  $|W|\le |S|e^{CP}$\\
  $|W\cap S|\ge c|S|$
};
\node[validationstage,right=of predict] (first) {
  \textbf{First validation}\\[5pt]
  Selected row in $W\cap S$\\[3pt]
  Test candidates in $U_T$\\
  Output $U_T'\subseteq U_T$\\
  $|U_T'\cap T|\ge .99|T|$
};
\node[validationstage,right=of first] (second) {
  \textbf{Second validation}\\[5pt]
  Selected row in $U_T'\cap T$\\[3pt]
  Test candidates in $U_S$\\
  Output $U_S'\subseteq U_S$\\
  $|U_S'\cap S|\ge .99|S|$
};
\draw[flow] (predict)--(first);
\draw[flow] (first)--(second);
\end{tikzpicture}
\caption{The description gives a first approximation; two separately
addressed validation tables give near-complete capture on both sides.
The second validation tests the original domain $U_S$, not $W$.
For either validation row, the constant-density comparison with iid
sampling holds at readiness and conditional on that row being produced,
not conditional on later success. No such comparison is asserted for
the description's accepted row.}
\label{geom:validation-figure}
\end{figure}


\begin{lemma}[The law of a produced test]\label{geom:produced-test}
Consider either test in Lemma~\ref{geom:validation}, conditional on
all data making that test ready, but not on any later event. Suppose its
hidden source $X_0$ is a subset of an original support $X$ with
$|X_0|\geq c|X|$. For every fixed opposite point $y$,
\begin{equation}\label{geom:exclusion-bound}
 \Prob\{\text{the test is produced and its ambient test excludes }y\}
 \leq Cq\frac{|X\cap N(y)|}{|X|}.
\end{equation}
Conditional on production, the row law is dominated by a constant times
the independent uniform row law on $X_0$. The bound
\eqref{geom:exclusion-bound} may be averaged over earlier readiness data
whenever its right-hand side is fixed.
\end{lemma}

\begin{proof}
Let $\nu$ be one proposal's law and $A$ its acceptance event. For
independent proposals and any deterministic cutoff $m$, the first
accepted row, conditional on there being one among the first $m$, has
law $\nu(\,\cdot\mid A)$. Indeed, the contribution of acceptance at
index $i$ is $(1-\nu(A))^{i-1}\nu(A)$ times that same conditional
law. Given that all entries lie in $X_0$, the proposal is an independent
uniform row on $X_0$, and the further acceptance probability is at least
$.9$. The claimed domination follows. Apply
\eqref{geom:pointwise-iid} and use
\[
 \frac{|X_0\cap N(y)|}{|X_0|}
 \leq c^{-1}\frac{|X\cap N(y)|}{|X|}.
\]
Multiplication by the probability of production can only decrease the
bound. No later success, including production of the other validation
test, has been conditioned on.
\end{proof}

\subsection{Reduction from higher projective dimensions}

\begin{proposition}[Repeated projection]\label{geom:projection}
For $j\geq4$, assume Lemma~\ref{geom:description} in projective
dimension $j-1$, for every fixed value of its hypothesis constant.
Then the lemma holds in dimension $j$.
\end{proposition}

\begin{proof}
For a point $z$ represented by a vector line $\langle z\rangle$, the
quotient vector space has dimension $j$ and hence projectivization
$\PG(j-1,q)$. Projection maps a point different from $z$ to its
image in this quotient; its fibers are the $q$ noncentral points of
lines through $z$. The dual operation retains hyperplanes containing
$z$, identified with hyperplanes of the quotient. Write $S_0,T_0$ for
the projected $S\setminus\{z\}$ and the cut $T\cap z^\perp$.
Apply the same operations to the decoder domains.

Lifting a quotient cap can enlarge it by a factor $q$, the size of a
projection fiber. We remove that enlargement by intersecting caps from
eight projections, always using the original $S,T$. At a given call let
$A\subseteq U_S$ be the intersection of all caps lifted so far, initially
$A=U_S$. The next center will preserve almost all of $S$ and limit
collisions among the points of this current set $A$; the latter property
makes successive intersections contract.
There exists a center, chosen before that
call's fresh tables, satisfying
\begin{equation}\label{geom:center-properties}
 \begin{aligned}
 |T_0|&\asymp |T|/q,&
 |U_T\cap z^\perp|&\leq C|U_T|/q,\\
 |S_0|&\geq(1-o(1))n,&
 e(S_0,T_0)&\leq C\tau |S_0||T_0|/q,\\
 \operatorname{coll}_z(A)&\leq C|A|^2/q^{j-1}.&&
 \end{aligned}
\end{equation}
Here $\operatorname{coll}_z(A)$ counts unordered distinct pairs of
$A\setminus\{z\}$ with the same projection.

To prove existence, choose $z$ uniformly in the ambient point space.
The variance identity gives
\[
 \Prob\left\{\left||T\cap z^\perp|-p_j|T|\right|
                >\frac{|T|}{3q}\right\}
 \leq Cq/|T|=o(1).
\]
Here $|T|\geq q^{(j+1)/2}e^{-b/2}\gg q$, and replacing the constant
$1/3$ by any other fixed positive constant is harmless. Markov's
inequality gives the stated cut-domain bound with arbitrarily high
fixed probability. A distinct pair projects together only if $z$ lies
on its joining line, an event of probability at most
$(q+1)/Q_j=O(q^{1-j})$. Thus
\[
 \E\operatorname{coll}_z(A)\leq C|A|^2/q^{j-1},\qquad
 \frac{\E\operatorname{coll}_z(S)}{n}
 \leq Cq^{(3-j)/2}=o(1).
\]
The latter implies projection loss $o(n)$ with probability $1-o(1)$:
each nonempty fiber loses at most its number of unordered pairs, and
omitting $z$ loses at most one more point. Finally, an original
incidence survives the dual cut with probability $p_j\leq C/q$.
Its expected count before collapsing fibers is at most
$C\tau n|T|/q^2$, and collapsing cannot increase it. A union bound,
using large fixed constants in the three Markov bounds, proves
\eqref{geom:center-properties}. This reasoning is valid for every
realized current $A$: fix that set and all earlier data, then combine
the two events of probability $1-o(1)$ with the three Markov bounds.
Their constants can be chosen so that the simultaneous event has
positive probability. Choose a center in that event before exposing
any of the fresh tables used by this call.

Transmit $z$ in $O(\sigma)$ bits and use a fixed coordinate convention
on the quotient. The new product is at least $q^je^{-b-O(1)}$ and the
density at most $C\tau/q$. The projected $S$ domain has size at most
$|U_S|$, whereas $|S_0|=(1-o(1))n$; on the other side both domain and
support shrink by comparable factors $q$. Hence both logarithmic gaps
increase by $O(1)$. These are the induction hypotheses, with changed
fixed constants. Reorient if necessary, apply induction, and then
Lemma~\ref{geom:validation}, which returns caps on both quotient
sides. In particular, the projected $S$ cap captures $.99|S_0|$ and
has size at most
\[
 \frac{Cq^j}{|T_0|}\leq \frac{Cq^{j+1}}{|T|}
 \leq m_0:=ne^{C_3P}.
\]
Lift this cap, excluding $z$, and intersect it with $A$.
The original support loss is at most $.01n$, plus the $o(n)$ total
fiber excess, plus one possible center: among the lost images there
are at most $.01|S_0|$ first preimages, and all additional preimages
are charged to $n-|S_0|-\boldsymbol 1_{\{z\in S\}}=o(n)$.
This bound concerns the original $S$, not merely the points retained
by earlier calls. Thus each call loses fewer than $.02n$ of its
original points, and losses from intersecting the lifted caps add.

The first call gives $|A_1|\leq qm_0$. For any later call let $k_i$
be the retained nonzero fiber sizes; there are at most $m_0$ of them.
Cauchy--Schwarz and \eqref{geom:center-properties} give
\[
 \begin{split}
 |A_{\mathrm{new}}|
 &\leq m_0+\sum_i(k_i-1)
 \leq m_0+\sqrt{m_0\sum_i k_i(k_i-1)}\\
 &\leq m_0+C|A|\sqrt{m_0/q^{j-1}}.
 \end{split}
\]
By \eqref{geom:product-upper}, the last coefficient is
\[
 \rho\leq e^{C_4P}q^{(3-j)/4}\leq q^{-1/4+o(1)}.
\]
After seven further calls,
$|A_8|\leq m_0(1+\rho+\cdots+\rho^6)+qm_0\rho^7=O(m_0)$,
while more than half of the original $S$ remains. Eight calls multiply
the cost and failure bounds by fixed constants. The center at each
stage depends only on data available before its fresh tables, so all
conditional failure bounds apply. The set $W=A_8$ proves the result.
\end{proof}

\subsection{Preparing dimensions two and three}

The remaining analysis will use a retained support $S'$ and, in one
case, an ordered list of planes. The list isolates planar concentrations;
for a query point we remove the training cell in its first containing
plane before measuring overlaps.

\begin{proposition}[Low-dimensional preparation]\label{geom:preparation}
Suppose \eqref{geom:sparse-hypotheses} holds with $j\in\{2,3\}$ and
$n=|S|\leq|T|$. Either enumeration already proves
Lemma~\ref{geom:description}, or the following preparation costs
$O(q+\sigma)$ bits. In dimension three there is also a possible
reduction to the dimension-two description lemma, with the same final
cost, size, capture, and failure bounds as
Lemma~\ref{geom:description}.

In the case not resolved by enumeration or that reduction, we have
$n>100qP$, a retained set $S'\subseteq S$ of size
$n'\in[n/2,n]$, and either an empty list or an ordered plane list
$\Pi_1,\ldots,\Pi_J$ in $\PG(3,q)$. In the latter case the sets
\begin{equation}\label{geom:cells}
 C_i=(S\cap\Pi_i)\setminus\bigcup_{h<i}\Pi_h
 \quad(1\leq i\leq J)
\end{equation}
partition $S'$, have nonincreasing sizes, and satisfy
$|C_i|<.04n'$. The decoder knows the list and all cell sizes, but need
not know $S'$ or the cells themselves.

Write
\begin{equation}\label{geom:lowdim-parameters}
 n=q^{1+v},\quad g=(v-1/2)\sigma,\quad
 \chi=P/100,\quad\chi_0=P/10000,\quad B=q^j/n'.
\end{equation}
For every ambient query point $x$, define $O_x$ to be the cell $C_i$
of its earliest containing listed plane, and $O_x=\varnothing$ if
there is none. Define
\begin{equation}\label{geom:outside-support}
 f_x=|O_x|/n'\leq .04,\qquad
 X_x=S'\setminus(O_x\cup\{x\}).
\end{equation}
If no list is present, $O_x=\varnothing$ for all $x$. When
$j=3$ and $g>\chi_0$, the preparation additionally ensures either
\begin{enumerate}[label=(\roman*)]
 \item every plane meets $S'$ in fewer than
 $K_p=n^{4/3}q^{-1}e^{-g/5}$ points and the list is empty; or
 \item the list is present and
 $J\leq n/K_p=q^{1/2}e^{-2g/15}$.
\end{enumerate}
\end{proposition}

\begin{proof}
If $n\leq100qP$, transmit $n$ and the index of $S$ among the
$n$-subsets of $U_S$. The cost is at most
\[
 O(\sigma)+\log_2\binom{|U_S|}{n}
 \leq C\bigl(\sigma+n(d_S+1)\bigr),
\]
which is within the description budget. It gives $W=S$ without
failure. Assume henceforth $n>100qP$, so $v>0$.

When $j=3$ and $g>\chi_0$, greedily choose a plane richest in the
remaining points of $S$, provided it contains at least $K_p$ of them,
and remove those points. Stop as soon as at least $n/2$ points have
been removed or no qualifying plane remains. In the first case retain
the captured prefix, with its list and cells; in the second retain the
residual without a list. Otherwise retain all of $S$ without a list.
The asserted size and cap properties follow directly from the stopping
rule. Greedy maximality makes the removed cell sizes nonincreasing.
Since each removal has size at least $K_p$, the list has length at
most $n/K_p=q^{1/2}e^{-2g/15}\leq q^{1/2}$. Plane identifiers and
cell sizes each cost $O(\sigma)$ bits, so the list costs $O(q)$.

Suppose now that a retained list has a cell of size at least $.04n'$,
and let $\Pi$ be its containing plane. Then
$S_1=S\cap\Pi$ has size at least $.02n$. The original sparsity and
Markov's inequality imply that at least half of $T$ consists of
hyperplanes having incidence proportion at most $C\tau/q$ into
$S_1$. For large $q$ this is less than one, so their restrictions to
$\Pi$ are nonzero projective covectors. Each nonzero projective
restriction has exactly $q$ projective lifts: normalize its restriction
to a fixed representative and choose the coefficient on a complement
of the vector space underlying $\Pi$.

Group the retained lifts by their restriction multiplicities. One
dyadic class $[h,2h)$, with $1\leq h\leq q$, carries at least
$c|T|/\sigma$ lifts. If $T_1$ is its set of distinct restrictions,
then
\[
 |T_1|\geq c|T|/(h\sigma),\qquad
 |S_1||T_1|\geq q^3e^{-b-O(\log\sigma)}.
\]
Every member of $T_1$ has incidence proportion at most $C\tau/q$
into $S_1$, since all its lifts have identical incidences there.
Transmit $\Pi$ and $h$. Use decoder domains $U_S\cap\Pi$ and the
nonzero restrictions having at least $h$ lifts in $U_T$. The latter
has size at most $|U_T|/h$, so its logarithmic gap increases by
$O(\log\sigma)$; the former gap increases by $O(1)$.

Apply the dimension-two description, with reorientation if needed,
and then Lemma~\ref{geom:validation}. The resulting cap on $S_1$
has size at most
\[
 \frac{Cq^3}{|T_1|}
 \leq\frac{Cq^4\sigma}{|T|}
 \leq Cn\sigma e^b\leq ne^{CP}
\]
and captures $.99|S_1|\geq cn$ points of $S$. Its cost is within
the required budget because $\log\sigma=o(P)$; its failure
probability has the required order. This deals with every large cell.
All other cases have the stated small-cell preparation. Notice that
$O_x$ is a subset of the training support: even when $x\notin S'$,
the earliest containing plane still specifies the same whole cell.
\end{proof}

\subsection{A polynomial bound for rich lines}

We give the needed algebraic estimate with its richness and
characteristic hypotheses explicit. It is used only when the line
richness substantially exceeds the cubic-root scale and planar
concentration cannot account for many rich lines. The interpolation
and tangent-plane method has antecedents in
\cite{GuthKatz2010,ElekesKaplanSharir2011}. Positive-characteristic
obstructions are discussed in \cite{EllenbergHablicsek2016}; here the
argument keeps every relevant surface degree below the characteristic.

\begin{lemma}[Rich lines with a plane cap]\label{geom:rich-lines}
Let $q$ run through sufficiently large primes. Let $\mathcal S$ be
at most $n\leq Cq^2$ points of $\PG(3,q)$, with at most $K_p$
points in every projective plane. Let $M>0$ satisfy
\begin{equation}\label{geom:rich-regime}
 \frac{M^3}{n}\longrightarrow\infty,\qquad
 \frac{M^2}{K_p}\longrightarrow\infty,\qquad
 cq\leq\frac nM\leq Cq^{41/40}.
\end{equation}
Then at most $Cn/M$ projective lines contain at least $M$ points of
$\mathcal S$. The constant is uniform once the displayed limits hold
uniformly in the family under consideration.
\end{lemma}

\begin{proof}
Work over the algebraic closure $\mathbb K$ of $\mathbb F_q$.
Choose an affine chart containing all points of $\mathcal S$, which
is possible because $\mathbb K$ is infinite. The plane cap remains
valid for planes over $\mathbb K$: rational points on such a plane
have rational vector-span dimension at most three. Otherwise four
linearly independent rational representatives would remain independent
over $\mathbb K$. The points therefore lie in an $\mathbb F_q$-plane
and obey the same cap.

Independently sample each point with probability
$A\sqrt{n/M^3}=o(1)$, where $A$ is a sufficiently large fixed
constant. The expected sample count on any $M$-rich line is at least
$A\sqrt{n/M}\geq cA\sqrt q$. There are at most $n^2$ lines to
test, since $M\to\infty$ and two distinct points determine a line.
Binomial tails and a union bound show that, simultaneously with
positive probability, every rich line contains more than
$(A/2)\sqrt{n/M}$ sampled points, and the total sample has size at
most $2A(n/M)^{3/2}$. The latter assertion follows from an upper
binomial tail, or from stochastic domination if $|\mathcal S|<n$.

There are $\binom{D+3}{3}$ monomials of degree at most $D$ in three
variables. Solving the homogeneous linear system requiring vanishing
at the sample therefore gives a nonzero polynomial of degree at most
$CA^{1/3}\sqrt{n/M}$. Choose $A$ large enough that this degree is
smaller than the guaranteed number of sampled points on each rich
line. Restriction to a line is a univariate polynomial, so it vanishes
identically there. Replace the polynomial by its squarefree part.
Its degree $D_0$ satisfies
\begin{equation}\label{geom:interpolating-degree}
 D_0=O(\sqrt{n/M})=o(M),\qquad
 D_0\leq Cq^{41/80}<q.
\end{equation}

We now count three groups of rich lines. Plane factors are controlled
by the plane cap. Among the other lines, those with few points incident
to three lines are controlled by the total point count. The remaining
lines will lie in the common zero set of the nonplane surface and an
auxiliary polynomial coprime to it, whose degrees bound their number.

First count rich lines lying in plane factors. There are at most
$\binom{D_0}{2}$ lines lying in two distinct such planes. For one
plane let $u$ be the number of remaining rich lines it contains.
Select $M'=\lfloor M\rfloor$ points of $\mathcal S$ on each.
If $r_a$ is the number of selected occurrences at point $a$, then
\[
 \sum_a r_a=uM',\qquad
 \sum_a r_a^2\leq uM'+u(u-1).
\]
The second inequality uses the fact that two distinct lines meet in
at most one point. Cauchy--Schwarz and the plane cap give
\[
 (uM')^2\leq K_p\bigl(uM'+u(u-1)\bigr).
\]
For $u>0$ this implies
$u\leq K_pM'/((M')^2-K_p)=o(M')$. Remove from each selected line
its intersections with all other plane factors, losing at most $D_0$
points. The remaining union in this plane has size at least
\[
 u(M'-D_0)-\binom u2\geq uM'/2.
\]
Such unions for different planes are disjoint. Thus the plane factors
account for $O(n/M+D_0^2)=O(n/M)$ rich lines.

Let $F_0$ be the squarefree product of the nonplane irreducible factors.
Every as yet uncounted rich line lies on $F_0$: the coordinate ring of
a line is an integral domain, so vanishing of a product on the line
implies vanishing of one factor. If there are no such factors we are
done. For coordinate vectors $e_1,e_2,e_3$, form
\begin{equation}\label{geom:flatness-polynomials}
 Q_{ij}=(e_i\times\nabla F_0)^{\mathsf T}
       \operatorname{Hess}(F_0)(e_j\times\nabla F_0),
 \qquad 1\leq i,j\leq3.
\end{equation}
At a point on at least three of the remaining lines, all $Q_{ij}$
vanish. At a singular point this follows from the gradient factors.
At a nonsingular point, the three distinct line directions lie in the
two-dimensional tangent vector space. Restricting $F_0$ to each line
shows that its Hessian quadratic form vanishes on each direction. A
binary quadratic vanishing on three distinct projective directions is
zero; since the characteristic is odd, its associated bilinear form
is also zero. The vectors $e_i\times\nabla F_0$ span that tangent
space, proving the assertion.

Among incidences with points of $\mathcal S$, points on at most two
remaining lines account for at most $2n$ incidences. Hence all but
$O(n/M)$ remaining rich lines have at least $M/2$ points incident
with three or more remaining lines. Since
$\deg Q_{ij}<3D_0=o(M)$, every such line lies on every $Q_{ij}$.

To obtain a polynomial coprime to $F_0$, we must show that the
$Q_{ij}$ do not all vanish identically on any of its irreducible
components. The reason is local: their vanishing makes all second
derivatives of a smooth formal graph zero, while the degree bound
below the prime characteristic forces a plane factor.

We claim that no irreducible factor $F$ of $F_0$ divides all these
polynomials. Its degree is less than the prime characteristic $q$,
so a nonconstant $F$ has a nonzero partial derivative. Rename its
variable $z$, and write the other variables as $x_1,x_2$. By
squarefreeness, $F$ and $(F_0)_z$ are coprime in
$\mathbb K(x_1,x_2)[z]$. Indeed, modulo $F$ the derivative is
$F_z(F_0/F)$, and neither factor is divisible by $F$. Clearing a
univariate Bezout identity gives
\[
 A F+B(F_0)_z=a(x_1,x_2)\ne0
\]
with polynomial coefficients. Choose $x_1,x_2$ so that $a$ and
the leading coefficient of $F$ in $z$ are nonzero. A root in $z$
then gives a point on $F$ where $(F_0)_z\ne0$.

Translate this point to the origin. The invertible coefficient of
$z$ in $F_0$ determines, recursively by total degree, a unique formal
series $z=h(x_1,x_2)$ with $F_0(x_1,x_2,h)=0$. The other factors
are units on this branch. If $F$ divided every $Q_{ij}$, those
polynomials would vanish after this substitution. To see the tangent
space statement over the formal series ring explicitly, write
$g=\nabla F_0(x_1,x_2,h)$. Its third coordinate $g_3$ is a unit,
and the kernel of the row vector $g$ has basis
\[
 t_1=(1,0,-g_1/g_3),\qquad t_2=(0,1,-g_2/g_3).
\]
Moreover $e_2\times g=g_3t_1$ and $e_1\times g=-g_3t_2$.
The vanishing $Q_{ij}$ therefore makes the Hessian bilinear form
zero on this entire formal tangent module. Twice differentiating
$F_0(x_1,x_2,h)=0$ gives
\[
 (F_0)_z h_{ab}=0\quad (a,b\in\{1,2\}),
\]
so all second partial derivatives of $h$ vanish. For a monomial
$x_1^u x_2^v$ of total degree less than $q$, these conditions require
\[
 u(u-1)=uv=v(v-1)=0\pmod q.
\]
Only the constant and linear monomials satisfy these conditions.
Writing $h_1$ for the affine-linear part, we obtain
$h-h_1\in(x_1,x_2)^q$. Thus the polynomial
$F_0(x_1,x_2,h_1)$ vanishes to order at least $q$. Its degree is
less than $q$, so it is identically zero. The plane $z=h_1$ would
be a factor of $F_0$, a contradiction. This proves the claim.

Because $\mathbb K$ is infinite, a constant linear combination $Q$
of the $Q_{ij}$ can be chosen coprime to $F_0$: for each of its
finitely many irreducible factors, combinations divisible by that
factor form a proper linear subspace of the coefficient space. We
finish by justifying the elementary common-line bound used here.

If coprime polynomials in three variables have degrees $D_1,D_2$,
they contain at most $2D_1D_2$ common lines. To see this, take any
finite collection of such lines and choose coordinates $(t,u,v)$
with $t$ nonconstant on each. Over the fraction fields in the other
two variables, univariate coprimality gives cleared identities
\[
 A F+B G=a(t,u)\ne0,\qquad
 C F+D G=b(t,v)\ne0.
\]
For a generic $c\in\mathbb K$, both right-hand sides remain
nonzero polynomials after $t=c$. Any common factor of the sliced
polynomials would divide both a polynomial in $u$ alone and one in
$v$ alone, and hence would be constant. Choose the slice also to
meet the selected lines in distinct points; only finitely many
values are forbidden by pairwise intersections of lines.
Choose coordinates on this slice so that these points have distinct
first coordinates. The resultant eliminating the second coordinate
is nonzero by coprimality and has degree at most $2D_1D_2$ in the
first coordinate. Indeed its Sylvester determinant has at most
$D_2$ columns with coefficients of degree at most $D_1$, and at
most $D_1$ columns with coefficients of degree at most $D_2$.
It vanishes at each of the distinct first coordinates, bounding
the collection. If both sliced polynomials omit the second
coordinate, their univariate coprimality gives no common point.
If just one omits it, say it is $f(u)$ and the other has degree
$m>0$ in the second coordinate, the resultant is $f(u)^m$;
it is again nonzero and vanishes at every common point.
As this bounds every finite collection, it bounds all common lines.

Apply the bound to $F_0,Q$, whose degrees are at most $D_0$ and
$3D_0$. The remaining lines number $O(D_0^2)=O(n/M)$. Together
with the earlier exceptional lines this proves the lemma.
\end{proof}

\subsection{Radial lines and overlap graphs}

The next two lemmas quantify the benefit of removing $O_x$. Planes
through a query point can overlap along a line, but after this removal
there are few lines carrying substantial retained mass, except at a
small set of query points.

\begin{lemma}[Radial rich lines]\label{geom:radial}
Consider the dimension-three output of
Proposition~\ref{geom:preparation}. A line $\ell$ through a query
point $x$ has strength
\[
 a_\ell(x)=q|X_x\cap\ell|/n'.
\]
Outside one exceptional set of $o(n)$ ambient points, simultaneously
for dyadic $a\in[q^{-10},2]$,
\begin{equation}\label{geom:radial-count}
 \#\{\ell\ni x:a_\ell(x)\geq a\}
 \leq q^{2-2v}e^\chi a^{-100}.
\end{equation}
In addition, whenever $g>\chi_0$ and $a>e^{-g/20}$ with
$0<a\leq2$, the total number of ambient center--line pairs
$(x,\ell)$ having strength at least $a$ is at most $Cq^2/a^2$.
\end{lemma}

\begin{proof}
The radial lines partition $X_x$, so their total strength is at most
$q$. Thus there are at most $q/a$ qualifying lines through any
center. Comparing this with \eqref{geom:radial-count} reduces to
\[
 e^{2g}a^{99}\leq e^\chi.
\]
This holds for $g\leq\chi_0$ and large $q$, since $a\leq2$ and
$2\chi_0<\chi$. It also holds when $g>\chi_0$ but
$a\leq e^{-g/20}$.

In the remaining regime put $M=n'a/q$. Since $n'\geq n/2$ and
$n=q^{3/2}e^g$,
\[
 \frac M{n^{1/3}}\geq c e^{(2/3-1/20)g},\qquad
 \frac{M^2}{K_p}\geq c e^{(2/3+1/10)g}.
\]
Also $cq\leq n/M\leq Cqe^{g/20}$, and
\eqref{geom:product-upper} gives $g\leq\sigma/2+O(1)$.
Hence all hypotheses of Lemma~\ref{geom:rich-lines} hold uniformly.
In the capped residual case, that lemma gives $Cn/M\leq Cq/a$
rich lines; each has $q+1$ rational centers. This yields at most
$Cq^2/a\leq Cq^2/a^2$ relevant pairs.

Suppose instead that the plane list is retained. Here own-cell removal
will force the center of a surviving rich line into a plane earlier
than the first listed plane containing that line. This gives a bound
on its possible centers. First, the list length satisfies
\[
 \frac JM\leq C e^{-(1+2/15-1/20)g}=o(1).
\]
A relevant line must be contained in a listed plane, since otherwise
it meets the union of the list in at most $J<M$ points. Let $i$
be its first containing plane. Earlier planes contribute at most
$J$ points to this line and later cells contribute none, so
$|C_i\cap\ell|\geq M-J\geq M/2$. Its center must be in an
earlier plane: otherwise $O_x=C_i$, leaving at most $J<M$ points
on this line in $X_x$. Since the line is not contained in any
earlier plane, there are at most $i-1$ possible centers for it.

Put $f_i=|C_i|/n'$. Fix a sufficiently large constant $C_0$.
If $a>C_0f_i$, the mean number of $C_i$ points on a line of
$\Pi_i$ is at most $2f_i n'/q$, much less than $M/2$.
The projective-plane incidence variance bounds the number of lines
with at least $M/2$ points by
\[
 \frac{Cq|C_i|}{M^2}
 \leq Cq^{2-v}f_i/a^2.
\]
Multiplying by at most $i$ centers and summing gives
\[
 Cq^{2-v}a^{-2}\sum_i i f_i
 \leq Cq^{2-v}J/a^2\leq Cq^2/a^2,
\]
because $\sum_i f_i=1$ and $J\leq q^v$. In the other case
$a\leq C_0f_i$, nonincreasing cell fractions and their unit sum
imply $i\leq C_0/a$; there are at most $C_0/a$ such indices.
Each corresponding plane has $O(q^2)$ lines and each line at
most $i$ possible centers. Their total is again $O(q^2/a^2)$.
This proves the center--line bound in all remaining cases.

For a fixed dyadic threshold, divide that bound by the right-hand
side of \eqref{geom:radial-count}. The number of centers where
the latter fails, divided by $n=q^{1+v}$, is at most
\[
 Cq^{v-1}e^{-\chi}a^{98}.
\]
Since $v\leq1+O(1/\sigma)$ and the sum of $a^{98}$ over all
dyadic $a\leq2$ is bounded, the union of these exceptional sets
has size $o(n)$, as asserted.
\end{proof}

\begin{lemma}[Overlap bounds for pencils]\label{geom:overlap}
Use the prepared data of Proposition~\ref{geom:preparation} in
dimension $j=2$ or $3$. For a query point $x$, let $\mathcal H_x$
be any family of hyperplanes through $x$ such that
\[
 \lambda'_H:=q|X_x\cap H|/n'\leq2\quad(H\in\mathcal H_x).
\]
For distinct $H,H'\in\mathcal H_x$ put
$\theta_{HH'}=q|X_x\cap H\cap H'|/n'$.
Let $E_a$ be the number of ordered pairs with $\theta_{HH'}\geq a$
and let $\Delta_a$ be the maximum degree of this graph. Outside
$o(n)$ ambient centers, simultaneously for dyadic
$a\in[q^{-10},2]$ and for every such family,
\begin{equation}\label{geom:overlap-bounds}
 E_a\leq B^2e^{2\chi}a^{-200},\qquad
 \Delta_a\leq Be^{2\chi}a^{-200}.
\end{equation}
Let $p$ be any even integer in
$[10000\sigma/P,10000\sigma/P+2]$ and put
$t_*=1/(100p)$. Outside a further set of at most $ne^{CP}$ ambient
centers, also
\begin{equation}\label{geom:strong-degree}
 \Delta_{t_*}\leq Be^{-3P}.
\end{equation}
The further exceptional set is used only for ambient output-size
control, not for support capture.
\end{lemma}

\begin{proof}
In dimension two, two distinct hyperplanes in the pencil are lines
intersecting only at $x$. Because $x\notin X_x$, all overlaps are
zero and all claims follow. Consider dimension three. The
intersection of two distinct planes through $x$ is a radial line;
each radial line belongs to $q+1$ planes and therefore accounts for
at most $Cq^2$ ordered pairs. Lemma~\ref{geom:radial} gives
\[
 E_a\leq Cq^{4-2v}e^\chi a^{-100}
 \leq B^2e^{2\chi}a^{-200}.
\]
The last inequality uses $B=q^3/n'\asymp q^{2-v}$ and $a\leq2$;
the factor $e^\chi$ absorbs fixed constants. Within one plane,
the radial lines partition its points other than $x$, so the sum
of their strengths is at most $2$. There are at most $2/a$
qualifying lines, each giving at most $q$ other planes. Thus
\begin{equation}\label{geom:crude-degree}
 \Delta_a\leq Cq/a.
\end{equation}
Since $B\geq cq$, this implies the second bound in
\eqref{geom:overlap-bounds}.

For the stronger assertion, $\log(1/t_*)=O(\log\sigma)=o(P)$.
If $(v-1)\sigma<-4P$, then
\[
 Be^{-3P}\geq cq e^P\geq Cq/t_*
\]
for large $q$, so \eqref{geom:crude-degree} suffices at every
center. Otherwise $n\geq q^2e^{-4P}$ and
$g\geq\sigma/2-4P>\chi_0$. Moreover
$t_*\gg e^{-g/20}$. The center--line estimate of
Lemma~\ref{geom:radial}, which does not require a dyadic threshold,
bounds the number of centers lying on even one radial line of
strength at least $t_*$ by
\[
 Cq^2/t_*^2\leq ne^{CP}.
\]
Remove these centers. At every remaining center
$\Delta_{t_*}=0$, proving \eqref{geom:strong-degree}.
\end{proof}

The geometry is now ready for the learning step. Its training support
is $S'$, the ordered plane data are known to the decoder, and the
own-cell rule in \eqref{geom:outside-support} is available for every
query point. Removing up to half of $S$ changes the sparse-incidence
bound by at most a factor of two. The overlap estimates are uniform
over all allowed pencil subfamilies, so those families may be selected
later using the prepared supports. Proposition~\ref{learn:short-row}
uses precisely these inputs to produce the cap required by
Lemma~\ref{geom:description}; the stronger ambient exceptional set
is never charged against support capture.
